\Needspace{4\baselineskip}
\begin{Example}
\label{ex:ell-two-hilbert-space}
Let $\ell^2$ be the set of all complex-valued sequences $x_n$ such that
\begin{equation}\label{eq:ell-two-sequence-square-norm}
\|x\|^2:=\sum_{n=1}^\infty|x_n|^2<\infty.
\end{equation}
The inner product is defined by
\begin{equation}\label{eq:ell-two-sequence-inner-product}
(x,y)=\sum_{n=1}^\infty\overline x_ny_n.
\end{equation}
It is easy to show that $\ell^2$ is an inner product space. We shall show it is complete. Let
\begin{equation}\label{eq:ell-two-cauchy-sequence-components}
u^{(n)}=(u_1^{(n)},u_2^{(n)},\ldots)
\end{equation}
be a Cauchy sequence. Then for any $\epsilon>0$, there is an $N(\epsilon)$ such that
\begin{equation}\label{eq:ell-two-cauchy-sequence-tail-bound}
\|u^{(n)}-u^{(m)}\|=\sqrt{\sum_{j=1}^\infty|u_j^{(n)}-u_j^{(m)}|^2}<\epsilon
\end{equation}
for $m,n>N(\epsilon)$. This implies in particular that $u_j^{(n)}$ is a Cauchy sequence for every fixed $j$. Let
\begin{equation}\label{eq:ell-two-coordinatewise-limit}
u_j=\lim_{n\to\infty}u_j^{(n)}.
\end{equation}
From \eqref{eq:ell-two-cauchy-sequence-tail-bound}, it follows that
\begin{equation}\label{eq:ell-two-finite-partial-sum-cauchy-bound}
\sum_{j=1}^k|u_j^{(n)}-u_j^{(m)}|^2\leq\epsilon^2
\end{equation}
for every $n,m>N(\epsilon)$ and every $k\in\mathbb{N}$. We let $m\to\infty$ and obtain
\begin{equation}\label{eq:ell-two-finite-partial-sum-limit-bound}
\sum_{j=1}^k|u_j^{(n)}-u_j|^2<\epsilon^2
\end{equation}
for $n>N(\epsilon)$, $k\in\mathbb{N}$. We now let $k\to\infty$ and conclude that $u^{(n)}\to u$ in $H$.
\end{Example}

\begin{Example}
\label{ex:lp-two-hilbert-space}
The space $L^2(\Omega)$ defined in Example~\ref{ex:continuous-functions-lp-norm}, with the inner product
\begin{equation}\label{eq:lp-two-inner-product}
(f,g)=\int_\Omega\overline{f(\bm{x})}g(\bm{x})\,\mathrm{d}\bm{x}
\end{equation}
is a Hilbert space. Here the integral in \eqref{eq:lp-two-inner-product} is defined in the sense of Remark~\ref{rem:lp-point-values-and-measure-zero}.
\end{Example}

\begin{Definition}
\label{def:separable-space}
A Hilbert space (or, more generally, a Banach space) is called \textbf{separable} if it contains a countable, dense subset.
\end{Definition}

Most spaces arising in applications are separable. Separability is important for the practical solution of problems, say, by discretization, because only countably many (well, in the real world, only finitely many) elements of the space can be represented in such a fashion. It is easy to see that $\ell^2$ is separable, because terminating sequences are dense. The space $L^2(\Omega)$ is also separable; see Problem~\ref{ex:lp-two-separable}.

\begin{Definition}
\label{def:hilbert-orthogonal-complement}
Let $H$ be a Hilbert space. We say that two elements of $H$, $x$ and $y$ are \textbf{orthogonal} if $(x,y)=0$. For any subspace $M$ of $H$, we define the \textbf{orthogonal complement} by
\begin{equation}\label{eq:hilbert-subspace-orthogonal-complement}
M^\perp=\{x\in H\mid(x,y)=0\ \forall y\in M\}.
\end{equation}

\end{Definition}

It is clear that $M^\perp$ is a closed subspace. If $M$ is also closed, then $H$ is the direct sum of $M$ and $M^\perp$: $H=M\oplus M^\perp$.

\begin{Theorem}{Projection theorem}
\label{thm:hilbert-orthogonal-projection}
Let $H$ be a Hilbert space and let $M$ be a closed subspace of $H$. Then every $u\in H$ has a unique decomposition $u=v+w$, where $v\in M$ and $w\in M^\perp$.
\end{Theorem}

\begin{Proof}
From elementary geometry, we expect $v$ to be the point in $M$ that is closest to $u$. Let us assume $u\notin M$ and let
\begin{equation}\label{eq:hilbert-projection-infimum-distance}
d:=\inf_{v^{\prime}\in M}\|u-v^{\prime}\|^2.
\end{equation}
Then there is a sequence $v_n\in M$ such that $d_n:=\|u-v_n\|^2$ converges to $d$. We shall prove that $v_n$ is a Cauchy sequence and take $v$ to be its limit.

Let $y$ be an arbitrary element of $M$ and let $\lambda$ be a scalar. Then $v_n+\lambda y\in M$, and hence
\begin{equation}\label{eq:projection-minimizing-sequence-perturbation}
\begin{aligned}
d&\leq\|u-(v_n+\lambda y)\|^2=\|u-v_n\|^2-\overline\lambda(y,u-v_n)\\
&\qquad-\lambda(u-v_n,y)+|\lambda|^2\|y\|^2.
\end{aligned}
\end{equation}
Setting $\lambda=(y,u-v_n)/\|y\|^2$, we conclude
\begin{equation}\label{eq:projection-minimizing-sequence-optimized-bound}
d\leq\|u-v_n\|^2-\dfrac{|(u-v_n,y)|^2}{\|y\|^2},
\end{equation}
which can be rewritten as
\begin{equation}\label{eq:projection-minimizing-sequence-pairing-bound}
|(u-v_n,y)|\leq\|y\|\sqrt{d_n-d}.
\end{equation}
Next, we observe that
\begin{equation}\label{eq:projection-minimizing-sequence-difference-pairing}
\begin{aligned}
|(v_m-v_n,y)|&\leq|(u-v_n,y)|+|(u-v_m,y)|\\
&\leq\|y\|(\sqrt{d_n-d}+\sqrt{d_m-d}).
\end{aligned}
\end{equation}
By setting $y=v_n-v_m$, we conclude that
\begin{equation}\label{eq:projection-minimizing-sequence-cauchy-bound}
\|v_n-v_m\|\leq\sqrt{d_n-d}+\sqrt{d_m-d},
\end{equation}
i.e., $v_n$ is a Cauchy sequence. Let $v$ be its limit. By taking limits in \eqref{eq:projection-minimizing-sequence-pairing-bound}, we find $(u-v,y)=0$ for every $y\in M$, i.e., $u-v\in M^\perp$. This proves the existence of the desired decomposition.

Suppose we had two decompositions $u=v+w=v^{\prime}+w^{\prime}$, then $v-v^{\prime}=w^{\prime}-w$, and hence
\begin{equation}\label{eq:orthogonal-projection-decomposition-uniqueness}
\|v-v^{\prime}\|^2=(v-v^{\prime},v-v^{\prime})=(v-v^{\prime},w^{\prime}-w)=0,
\end{equation}
since $v-v^{\prime}\in M$ and $w^{\prime}-w\in M^\perp$.
\end{Proof}

\begin{Corollary}
\label{cor:hilbert-density-orthogonal-complement}
Let $H$ be a Hilbert space. A subspace $M$ of $H$ is dense iff $M^\perp=\{0\}$.
\end{Corollary}

\begin{Proof}
Suppose $v\in M^\perp$. Then $(v,u)=0$ for every $u\in M$. If $M$ is dense, we conclude that $(v,u)=0$ for every $u\in H$ by taking limits. Setting $u=v$ yields $v=0$.

On the other hand, the closure of $M$ has an orthogonal complement according to the preceding theorem. If $M$ is not dense, then this orthogonal complement must contain nonzero vectors.
\end{Proof}
